\documentclass[11pt]{article}

% Shared notes settings for Elements of AIML lectures (UPES).
% Keep this file free of \documentclass to allow reuse via \input.

\usepackage[utf8]{inputenc}
\usepackage[T1]{fontenc}
\usepackage[a4paper,margin=1in]{geometry}
\usepackage{hyperref}
\usepackage{xurl}
\usepackage{booktabs}
\usepackage{amsmath}
\usepackage{amssymb}
\usepackage{listings}
\usepackage{xcolor}
\usepackage{enumitem}
\usepackage{parskip}

% Code listings (general-purpose).
\lstset{
  basicstyle=\ttfamily\small,
  keywordstyle=\color{blue},
  commentstyle=\color{gray},
  stringstyle=\color{teal},
  showstringspaces=false,
  columns=fullflexible,
  frame=single,
  framerule=0.2pt,
  breaklines=true
}

\setlist[itemize]{topsep=0.3em,itemsep=0.2em}
\setlist[enumerate]{topsep=0.3em,itemsep=0.2em}

% Simple per-section source line.
\newcommand{\notesources}[1]{\par\noindent{\small\color{gray}\textit{Sources:} \sloppy #1\par}}

% Unit-wise source strings for this subject (used by slides/notes).
% Path is relative to the lecture `latex/` folder (the compilation CWD).
\IfFileExists{../../../_shared/aiml-sources.tex}{%
  % Source strings for Elements of AIML (used by slides + notes).

\newcommand{\AIMLRepoURL}{https://github.com/tali7c/Elements-of-AIML}

\newcommand{\AIMLLabSources}{%
Provost and Fawcett, \textit{Data Science for Business}.;%
McKinney, \textit{Python for Data Analysis}.;%
WEKA Documentation (Waikato Environment for Knowledge Analysis), accessed 2026-02-28.;%
Matplotlib and Seaborn documentation, accessed 2026-02-28.%
}

\newcommand{\AIMLUnitISources}{%
Rich and Knight, \textit{Artificial Intelligence}, McGraw Hill, 2017.;%
Russell and Norvig, \textit{Artificial Intelligence: A Modern Approach} (4e), Ch. 1--2 (Introduction, Intelligent Agents).;%
Poole and Mackworth, \textit{Artificial Intelligence: Foundations of Computational Agents} (2e), selected sections.%
}

\newcommand{\AIMLUnitIISources}{%
Russell and Norvig, \textit{Artificial Intelligence: A Modern Approach} (4e), Ch. 7--9 (Logical Agents, First-Order Logic, Inference).;%
Huth and Ryan, \textit{Logic in Computer Science} (2e), selected sections on propositional and first-order logic.;%
Genesereth and Nilsson, \textit{Logical Foundations of Artificial Intelligence}, selected chapters.%
}

\newcommand{\AIMLUnitIIISources}{%
Mueller and Massaron, \textit{Machine Learning for Dummies}, For Dummies, 2016.;%
Mitchell, \textit{Machine Learning}, selected chapters on learning basics and evaluation.;%
Scikit-learn documentation, accessed 2026-02-28.%
}

\newcommand{\AIMLUnitIVSources}{%
Mueller and Massaron, \textit{Machine Learning for Dummies}, For Dummies, 2016.;%
James, Witten, Hastie, and Tibshirani, \textit{An Introduction to Statistical Learning} (2e), selected chapters.;%
Scikit-learn documentation, accessed 2026-02-28.%
}

\newcommand{\AIMLUnitVSources}{%
NIST, \textit{AI Risk Management Framework (AI RMF 1.0)}, 2023.;%
Barocas, Hardt, and Narayanan, \textit{Fairness and Machine Learning} (online book), selected sections.;%
Knaflic, \textit{Storytelling with Data}, selected chapters on communicating results.%
}
%
}{%
  % Faculty copies live under `private/` so their compile CWD is different.
  \IfFileExists{../../../../public/_shared/aiml-sources.tex}{%
    % Source strings for Elements of AIML (used by slides + notes).

\newcommand{\AIMLRepoURL}{https://github.com/tali7c/Elements-of-AIML}

\newcommand{\AIMLLabSources}{%
Provost and Fawcett, \textit{Data Science for Business}.;%
McKinney, \textit{Python for Data Analysis}.;%
WEKA Documentation (Waikato Environment for Knowledge Analysis), accessed 2026-02-28.;%
Matplotlib and Seaborn documentation, accessed 2026-02-28.%
}

\newcommand{\AIMLUnitISources}{%
Rich and Knight, \textit{Artificial Intelligence}, McGraw Hill, 2017.;%
Russell and Norvig, \textit{Artificial Intelligence: A Modern Approach} (4e), Ch. 1--2 (Introduction, Intelligent Agents).;%
Poole and Mackworth, \textit{Artificial Intelligence: Foundations of Computational Agents} (2e), selected sections.%
}

\newcommand{\AIMLUnitIISources}{%
Russell and Norvig, \textit{Artificial Intelligence: A Modern Approach} (4e), Ch. 7--9 (Logical Agents, First-Order Logic, Inference).;%
Huth and Ryan, \textit{Logic in Computer Science} (2e), selected sections on propositional and first-order logic.;%
Genesereth and Nilsson, \textit{Logical Foundations of Artificial Intelligence}, selected chapters.%
}

\newcommand{\AIMLUnitIIISources}{%
Mueller and Massaron, \textit{Machine Learning for Dummies}, For Dummies, 2016.;%
Mitchell, \textit{Machine Learning}, selected chapters on learning basics and evaluation.;%
Scikit-learn documentation, accessed 2026-02-28.%
}

\newcommand{\AIMLUnitIVSources}{%
Mueller and Massaron, \textit{Machine Learning for Dummies}, For Dummies, 2016.;%
James, Witten, Hastie, and Tibshirani, \textit{An Introduction to Statistical Learning} (2e), selected chapters.;%
Scikit-learn documentation, accessed 2026-02-28.%
}

\newcommand{\AIMLUnitVSources}{%
NIST, \textit{AI Risk Management Framework (AI RMF 1.0)}, 2023.;%
Barocas, Hardt, and Narayanan, \textit{Fairness and Machine Learning} (online book), selected sections.;%
Knaflic, \textit{Storytelling with Data}, selected chapters on communicating results.%
}
%
  }{%
    % Source strings for Elements of AIML (used by slides + notes).

\newcommand{\AIMLRepoURL}{https://github.com/tali7c/Elements-of-AIML}

\newcommand{\AIMLLabSources}{%
Provost and Fawcett, \textit{Data Science for Business}.;%
McKinney, \textit{Python for Data Analysis}.;%
WEKA Documentation (Waikato Environment for Knowledge Analysis), accessed 2026-02-28.;%
Matplotlib and Seaborn documentation, accessed 2026-02-28.%
}

\newcommand{\AIMLUnitISources}{%
Rich and Knight, \textit{Artificial Intelligence}, McGraw Hill, 2017.;%
Russell and Norvig, \textit{Artificial Intelligence: A Modern Approach} (4e), Ch. 1--2 (Introduction, Intelligent Agents).;%
Poole and Mackworth, \textit{Artificial Intelligence: Foundations of Computational Agents} (2e), selected sections.%
}

\newcommand{\AIMLUnitIISources}{%
Russell and Norvig, \textit{Artificial Intelligence: A Modern Approach} (4e), Ch. 7--9 (Logical Agents, First-Order Logic, Inference).;%
Huth and Ryan, \textit{Logic in Computer Science} (2e), selected sections on propositional and first-order logic.;%
Genesereth and Nilsson, \textit{Logical Foundations of Artificial Intelligence}, selected chapters.%
}

\newcommand{\AIMLUnitIIISources}{%
Mueller and Massaron, \textit{Machine Learning for Dummies}, For Dummies, 2016.;%
Mitchell, \textit{Machine Learning}, selected chapters on learning basics and evaluation.;%
Scikit-learn documentation, accessed 2026-02-28.%
}

\newcommand{\AIMLUnitIVSources}{%
Mueller and Massaron, \textit{Machine Learning for Dummies}, For Dummies, 2016.;%
James, Witten, Hastie, and Tibshirani, \textit{An Introduction to Statistical Learning} (2e), selected chapters.;%
Scikit-learn documentation, accessed 2026-02-28.%
}

\newcommand{\AIMLUnitVSources}{%
NIST, \textit{AI Risk Management Framework (AI RMF 1.0)}, 2023.;%
Barocas, Hardt, and Narayanan, \textit{Fairness and Machine Learning} (online book), selected sections.;%
Knaflic, \textit{Storytelling with Data}, selected chapters on communicating results.%
}
%
  }%
}


\title{Elements of AIML\\Unit 04 -- Lecture 03 Notes\\Clustering, Reinforcement Learning \& Semi-Supervised ML}
\author{Tofik Ali}
\date{\today}

\begin{document}
\maketitle
\tableofcontents

\section{Lecture Overview}
This lecture covers three important ML paradigms beyond basic supervised learning:
\begin{itemize}[leftmargin=*]
  \item \textbf{Clustering} (unsupervised),
  \item \textbf{Reinforcement learning} (interaction + reward),
  \item \textbf{Semi-supervised learning} (limited labels + many unlabeled examples).
\end{itemize}
\notesources{\AIMLUnitIVSources}

\section{How to use these notes (self-study)}
This lecture has three mini-topics. Study them in this order:
(1) clustering (because it uses only data $X$), (2) semi-supervised learning (because it still starts from labeled examples), and (3) reinforcement learning (because it changes the setting from datasets to interaction over time).
For each topic, focus on the \textbf{objective} and the \textbf{main pitfalls}.

\section{Learning Outcomes}
\begin{itemize}[leftmargin=*]
  \item Explain clustering and describe k-means steps and objective.
  \item Identify practical issues in clustering (choice of $k$, initialization, scaling).
  \item Explain key RL terminology and the RL learning goal.
  \item Describe semi-supervised learning and risks of pseudo-labeling.
\end{itemize}

\section{Keywords}
\begin{itemize}[leftmargin=*]
  \item \textbf{Cluster:} group of similar points.
  \item \textbf{k-means:} clustering algorithm minimizing within-cluster squared distances.
  \item \textbf{State/action/reward/policy:} core RL terms.
  \item \textbf{Semi-supervised:} learning using both labeled and unlabeled data.
\end{itemize}

\section{Abbreviations and symbols}
\begin{itemize}[leftmargin=*]
  \item $k$: number of clusters; $\mu_j$: centroid of cluster $j$; $c(i)$: cluster assignment of point $x_i$.
  \item Inertia/objective: $\sum_i \|x_i-\mu_{c(i)}\|^2$ (within-cluster sum of squares).
  \item RL: $s$ state, $a$ action, $r$ reward, $\pi$ policy, $\gamma$ discount factor.
  \item $Q(s,a)$: action-value (Q-function); $\alpha$: learning rate.
  \item Semi-supervised: labeled set $L$, unlabeled set $U$, pseudo-labels.
\end{itemize}

\section{Three paradigms at a glance}
\begin{itemize}[leftmargin=*]
  \item \textbf{Clustering:} no labels, no rewards. The question is ``what groups already exist in the data?''.
  \item \textbf{Semi-supervised learning:} a few labels plus many unlabeled examples. The question is ``how can unlabeled data help a supervised model?''.
  \item \textbf{Reinforcement learning:} no fixed label per row; instead the agent receives delayed rewards while interacting with an environment.
\end{itemize}
This comparison matters because the three paradigms use different feedback signals and therefore require different evaluation habits.

\section{Clustering}
\subsection{What is clustering?}
Clustering is an unsupervised task: we group data points to discover structure.
Unlike classification, there are no labels provided.

\subsection{Use-cases}
\begin{itemize}[leftmargin=*]
  \item customer segmentation,
  \item grouping similar documents/images,
  \item anomaly detection (points far from clusters).
\end{itemize}

\section{k-means}
\subsection{Algorithm}
k-means alternates between assignment and update:
\begin{enumerate}[leftmargin=*]
  \item initialize $k$ centers,
  \item assign each point to the nearest center,
  \item recompute each center as mean of assigned points,
  \item repeat until assignments stabilize.
\end{enumerate}

\subsection{Objective}
k-means minimizes within-cluster sum of squares:
\[
  \sum_{i=1}^{n}\left\|x_i - \mu_{c(i)}\right\|^2
\]

\subsection{Why the objective makes sense}
The objective encourages:
\begin{itemize}[leftmargin=*]
  \item \textbf{compactness:} points in a cluster should be close to the centroid,
  \item \textbf{separation:} different clusters represent different regions of the feature space.
\end{itemize}
Because it uses Euclidean distance, k-means works best when clusters are roughly spherical and similar in size.

\subsection{Practical notes}
\begin{itemize}[leftmargin=*]
  \item scaling matters (distance-based),
  \item initialization affects result; multiple runs help,
  \item choosing $k$ is not automatic (elbow method is a heuristic).
\end{itemize}

\subsection{Choosing $k$ and evaluating clustering}
Since clustering has no ground-truth labels, we use heuristics and domain checks:
\begin{itemize}[leftmargin=*]
  \item \textbf{Elbow method:} plot inertia vs $k$ and look for diminishing returns.
  \item \textbf{Silhouette score (intuition):} compares how close a point is to its own cluster vs other clusters.
  \item \textbf{Stability:} run k-means multiple times; stable clusters increase confidence.
  \item \textbf{Domain interpretability:} clusters should map to meaningful groups (e.g., customer segments).
\end{itemize}

\subsection{Mini worked example (1D)}
Points: $[1,2,8,9]$ with $k=2$.
Initialize $\mu_1=2$, $\mu_2=9$.
\begin{itemize}[leftmargin=*]
  \item Assign: $\{1,2\}\rightarrow \mu_1$ and $\{8,9\}\rightarrow \mu_2$.
  \item Update: $\mu_1=\text{mean}(1,2)=1.5$, $\mu_2=\text{mean}(8,9)=8.5$.
\end{itemize}
Next iteration keeps the same assignments, so the algorithm converges.

\section{Semi-supervised learning}
\subsection{Motivation}
Often labels are expensive (medical, legal), but unlabeled data is abundant.

\subsection{Self-training (one simple approach)}
\begin{enumerate}[leftmargin=*]
  \item train a model on labeled data,
  \item predict labels for unlabeled data,
  \item add high-confidence pseudo-labels to training,
  \item repeat.
\end{enumerate}
Risk: if pseudo-labels are wrong, errors can be reinforced.

\subsection{Pseudo-labeling failure example}
Imagine a medical-image dataset with very few positive cases.
If the first model is already biased toward the majority class, it may assign high-confidence negative pseudo-labels to borderline positive cases.
Those wrong pseudo-labels are then fed back into training, making the model even less sensitive to rare positives.
So semi-supervised learning must be validated carefully; ``more data'' is not automatically ``better data''.

\subsection{Other semi-supervised approaches (brief)}
\begin{itemize}[leftmargin=*]
  \item \textbf{Co-training}: two models teach each other using different feature ``views'' (when applicable).
  \item \textbf{Label propagation}: build a similarity graph and spread labels to nearby points.
  \item \textbf{Consistency regularization}: encourage stable predictions under perturbations (modern deep learning idea).
\end{itemize}

\section{Reinforcement learning (overview)}
\subsection{Core idea}
RL learns by interaction. At each time step:
\begin{itemize}[leftmargin=*]
  \item agent observes state $s$,
  \item chooses action $a$,
  \item receives reward $r$ and new state $s'$.
\end{itemize}
Goal: learn a policy $\pi(a|s)$ that maximizes cumulative reward.
This is where the agent concepts from \textbf{Unit 01 Lecture 02} become directly useful again.

\subsection{Return and discounting}
We usually care about the \textbf{return} (total future reward), not just immediate reward:
\[
  G_t = r_{t+1} + \gamma r_{t+2} + \gamma^2 r_{t+3} + \dots
\]
where $\gamma \in [0,1)$ is the \textbf{discount factor}.
\begin{itemize}[leftmargin=*]
  \item $\gamma$ close to 0: focus on short-term reward.
  \item $\gamma$ close to 1: focus on long-term reward.
\end{itemize}

\subsection{Q-learning (numeric intuition)}
One common RL idea is to update the value of an action using both the immediate reward and the estimated future value.
Suppose:
\begin{itemize}[leftmargin=*]
  \item old estimate $Q(s,a)=2.0$,
  \item immediate reward $r=1$,
  \item best next-state value $\max_{a'}Q(s',a')=3.0$,
  \item learning rate $\alpha=0.5$ and discount $\gamma=0.9$.
\end{itemize}
Then the target is:
\[
  r + \gamma \max_{a'}Q(s',a') = 1 + 0.9 \times 3 = 3.7
\]
and the updated value becomes:
\[
  Q_{\text{new}}(s,a) = 2.0 + 0.5(3.7 - 2.0) = 2.85
\]
Interpretation: the old estimate moves partway toward a better estimate based on reward plus promising future outcomes.

\subsection{Why RL is harder}
\begin{itemize}[leftmargin=*]
  \item feedback can be delayed,
  \item exploration vs exploitation trade-off,
  \item environment can be stochastic.
\end{itemize}

\section{Practice (with answers)}
\subsection*{P1. Why does scaling matter for k-means?}
\textbf{Answer:} k-means is distance-based. If one feature has a large numeric scale, it dominates the distance calculation and therefore dominates clustering.

\subsection*{P2. Give one RL example with state, action, and reward.}
\textbf{Answer:} robot navigation: state = robot position; action = move; reward = +1 for reaching goal, negative for hitting obstacles, small negative per step.

\subsection*{P3. What is the main risk of self-training?}
\textbf{Answer:} the model can add incorrect pseudo-labels and reinforce its own mistakes, reducing performance.

\section{Common pitfalls (and how to avoid them)}
\begin{itemize}[leftmargin=*]
  \item \textbf{Using k-means on unscaled data:} distance-based methods are dominated by large-scale features.
  \item \textbf{Assuming k-means finds the global optimum:} it can converge to local minima; use multiple initializations.
  \item \textbf{Confusing clustering with classification:} clusters are discovered structure, not ground-truth labels.
  \item \textbf{Ignoring exploration in RL:} without exploration the agent may never discover better actions.
  \item \textbf{Self-training error reinforcement:} pseudo-labeling can amplify early mistakes; validate carefully.
\end{itemize}

\section{Exit questions (with answers)}

\subsection*{Q1. What is the k-means objective in words?}
\textbf{Answer:} minimize the total squared distance between points and their assigned cluster centers.

\subsection*{Q2. Why can k-means produce different results on different runs?}
\textbf{Answer:} because random initialization can lead to different local minima; multiple initializations reduce this issue.

\subsection*{Q3. Define state, action, reward, policy.}
\textbf{Answer:} state = situation description; action = choice made by agent; reward = feedback signal; policy = mapping from states to actions (or action probabilities).

\subsection*{Q4. One benefit and one risk of semi-supervised learning?}
\textbf{Answer:} benefit: uses unlabeled data to improve performance; risk: wrong pseudo-labels can reinforce errors.

\section{Summary}
Clustering discovers structure without labels, RL learns from rewards, and semi-supervised learning leverages unlabeled data when labels are scarce.
\notesources{\AIMLUnitIVSources}

\end{document}
